\documentclass[11pt,a4paper]{article}
\usepackage[T1]{fontenc}
\usepackage[utf8]{inputenc}
\usepackage{newtxtext}
\usepackage[margin=27mm,headheight=15pt]{geometry}
\usepackage{amsmath,amsthm,mathtools}
\usepackage{newtxmath}
\usepackage{microtype}
\usepackage{booktabs,longtable,array}
\usepackage{enumitem}
\usepackage{needspace}
\usepackage{listings}
\usepackage{xcolor}
\usepackage{xurl}
\usepackage{fancyhdr}
\usepackage{hyperref}
\definecolor{linkink}{HTML}{35594D}
\hypersetup{colorlinks=true,linkcolor=linkink,citecolor=linkink,urlcolor=linkink,
 pdftitle={Radical Denesting: A Unified Literature and Implementation Report},
 pdfauthor={Unified research report prepared for Vladimir Reshetnikov},
 pdfsubject={Radical denesting, Galois theory, algorithms, computer algebra, annotated bibliography}}
\urlstyle{same}
\pagestyle{fancy}\fancyhf{}
\fancyhead[L]{\small Radical denesting}
\fancyhead[R]{\small Unified literature report}
\fancyfoot[C]{\thepage}
\renewcommand{\headrulewidth}{0.3pt}
\setlength{\parskip}{4pt}
\setlength{\parindent}{1.2em}
\setlength{\emergencystretch}{2em}
\setcounter{tocdepth}{2}
\numberwithin{equation}{section}
\newtheorem{theorem}{Theorem}[section]
\newtheorem{proposition}[theorem]{Proposition}
\newtheorem{corollary}[theorem]{Corollary}
\theoremstyle{definition}\newtheorem{definition}[theorem]{Definition}
\theoremstyle{remark}\newtheorem{remark}[theorem]{Remark}
\newcommand{\Q}{\mathbb Q}\newcommand{\R}{\mathbb R}\newcommand{\C}{\mathbb C}\newcommand{\Z}{\mathbb Z}
\DeclareMathOperator{\Gal}{Gal}\DeclareMathOperator{\sgn}{sgn}\DeclareMathOperator{\depth}{depth}
\newcommand{\doi}[1]{\href{https://doi.org/#1}{\nolinkurl{doi:#1}}}
\newcommand{\access}[2]{\href{#1}{#2}}
\newcommand{\status}[1]{\textsf{\footnotesize[#1]}}
\newcommand{\bibgroup}[1]{\par\Needspace{7\baselineskip}\bigskip\item[]\textbf{#1}\par\smallskip}
\lstset{basicstyle=\small\ttfamily,columns=fullflexible,keepspaces=true,
 breaklines=true,showstringspaces=false,frame=single,framerule=0.25pt,
 rulecolor=\color{black!30},xleftmargin=0.4em,xrightmargin=0.4em,
 aboveskip=8pt,belowskip=8pt,tabsize=4}
\begin{document}
\begin{titlepage}
\thispagestyle{empty}
\vspace*{18mm}
{\large\scshape Literature, theory, and computer algebra\par}
\vspace{10mm}
{\Huge\bfseries Radical Denesting\par}
\vspace{4mm}
{\LARGE A Unified Literature and\par Implementation Report\par}
\vspace{10mm}
{\large Classical criteria, field-theoretic algorithms,\par Ramanujan identities, and computational practice\par}
\vspace{15mm}
\noindent\rule{\textwidth}{0.6pt}
\vspace{5mm}
\noindent This report consolidates the four supplied research reports, removes repeated accounts and duplicate bibliographic records, corrects material discrepancies, and adds selected primary sources and reproducible checks. Conference papers and their journal expansions are identified as related versions rather than counted as independent discoveries.
\vspace{6mm}
\noindent The emphasis is finite algebraic denesting. Infinite nested radicals, polynomial solving, radical normalization, and identity testing are retained as related subjects, with their different questions made explicit.
\vfill
\noindent Prepared for Vladimir Reshetnikov\\
Research and software documentation checked: 5 September 2026
\end{titlepage}

\begin{abstract}
Radical denesting asks whether a number already represented by nested radicals has an equivalent expression with fewer radical layers. This survey unifies the classical square-root criteria, the algorithmic work of Caviness--Fateman, Borodin--Fagin--Hopcroft--Tompa, Zippel, Landau, Horng--Huang, and Bl\"omer, later results on Ramanujan-type identities and restricted higher-index radicals, and the interfaces actually offered by computer-algebra systems. It separates four easily confused measures: expression depth, algebraic degree, output size, and the size of a splitting-field representation. A branch-aware elementary development supplies proofs and diagnostic examples. An annotated bibliography links primary articles, theses, books, software documentation, and useful discussions. Substantive corrections include the alleged equivalence of rational cubic denesting to integer factorization, the misreading of SymPy's iteration bound as a depth bound, and several bibliographic and real-versus-complex ambiguities. Small exact computations accompany the survey; they are not presented as benchmarks of systems that were not executed.
\end{abstract}
\tableofcontents
\clearpage

\section{Scope, conventions, and a map of the literature}
\label{sec:scope}

\subsection{What is being minimized?}
A radical expression over a field $K$ is built from constants in $K$, arithmetic operations, and specified roots. Its syntactic depth is defined recursively by
\begin{align}
\depth(c)&=0 &&(c\in K),\\
\depth(E\mathbin\circ F)&=\max\{\depth(E),\depth(F)\}
 &&(\circ\in\{+,-,\times,/\}),\\
\depth(\sqrt[n]{E})&=1+\depth(E),
\end{align}
with nonzero denominators and a stated interpretation of the root. Integer powers need not introduce radical layers. For rational exponents, reducing the exponent and deciding which roots are primitive operations are part of the representation convention.

Denesting can mean removing one layer, reaching depth one, or achieving the smallest possible depth. These are different specifications. So are denesting with square roots only, allowing fourth roots, allowing arbitrary real radicals, and allowing complex radicals and roots of unity. The general literature must be read with the base field and admissible operations fixed; Landau's treatment makes the roots-of-unity convention particularly important \cite{landau92,landauNote,horng99,bloemer00}.

For example,
\begin{equation}\label{eq:standard}
\sqrt{5+2\sqrt6}=\sqrt2+\sqrt3
\end{equation}
lowers depth from two to one but does not lower algebraic degree. Both sides generate the same degree-four field. In contrast, writing an algebraic number as a \texttt{Root} object suppresses visible radical signs without producing a radical expression of smaller depth. Such a representation is excellent for exact arithmetic, but it answers a different question.

Throughout the elementary sections, square roots of nonnegative real numbers are nonnegative, and odd roots of real numbers are \emph{real} odd roots. This differs from the principal complex power $z^{1/n}$ used in many symbolic interfaces. The distinction will matter even for the innocent-looking identity
\begin{equation}\label{eq:fifth-real}
\sqrt[5]{41-29\sqrt2}=1-\sqrt2.
\end{equation}
The radicand and the right side are negative. Equation~\eqref{eq:fifth-real} is an identity of real fifth roots, not of principal complex fifth powers.

\subsection{Five neighboring questions}
\begin{description}[style=nextline,leftmargin=0pt,labelwidth=0pt]
\item[Denesting.] Find a shallower radical expression, or prove none exists in a specified class.
\item[Normalization and zero testing.] Put algebraic expressions into a form in which equality can be decided, possibly using a field basis rather than simple radicals.
\item[Solving polynomials by radicals.] Determine whether roots admit radical representations at all, and construct them when appropriate. Denesting starts with an already radical-solvable element.
\item[Degree reduction.] Express an element using generators of smaller algebraic degree. The radical depth of the resulting expression need not be smaller.
\item[Infinite-radical evaluation.] Prove convergence of a sequence of finite approximants and identify its limit. This is an analytic problem, not finite denesting.
\end{description}
These distinctions explain why RADCAN, Landau's denesting algorithm, Girstmair's degree-reduction results, and Herschfeld's convergence theorem belong in the same survey but cannot be substituted for one another \cite{caviness76,landau92,girstmair,herschfeld}.

\subsection{The main research strands}
The foundational algorithmic sequence runs from radical normalization in the 1970s through the two adjacent 1985 \emph{Journal of Symbolic Computation} papers of Borodin et al. and Zippel, then through Landau's general framework, Horng--Huang's minimum-depth constructions, and Bl\"omer's representation-sensitive algorithms. Jeffrey--Rich provides the practical bridge to recursive simplification \cite{caviness76,bfht,zippel85,landau92,horng90,bloemer92,bloemer00,jeffrey}.

A second strand begins with Ramanujan's identities and leads to units, pure cubic fields, quartic rational-root criteria, Cardan polynomials, and elliptic-curve constructions. Honsbeek, Sury, Krepkii--Pimenov, Antipov--Pimenov, Girstmair, Cavallo, Osipov--Kytmanov, and Tonien supply distinct pieces of this picture \cite{berndt97,berndt98,honsbeek05,sury,krepkii,antipov,girstmair,cavallo,osipov21,tonien}.

The computational-complexity strand studies succinct sums and circuits containing radical constants. Its progress does not automatically translate into an efficient formula-producing denester. Conversely, the failure of a heuristic denester to simplify an identity does not imply that equality is undecidable \cite{bloemer91,bloemer98,rit,gaillard}.

\subsection{How this consolidation treats evidence}
The supplied documents were discovery aids, not independent corroboration of a repeated claim. Mathematical results below are linked to their original sources, or proved directly when elementary. Bibliographic entries distinguish relevant text consulted (\status{T}), publisher or institutional metadata/abstract checked (\status{A}), official documentation or source inspected (\status{D}), and retained background or bibliographic leads not audited at theorem level (\status{B}). A label does not certify an entire paper or software package.

No claim is made that every linked paper was read in full, that the bibliography is exhaustive, or that all listed systems were executed. The exact test run in Section~\ref{sec:tests} used SymPy 1.14.0. The 2026 Osipov addition is based on the publisher's abstract because the full text required authenticated access. Appendix~\ref{sec:audit} records corrections and unresolved leads.

\section{Square roots: exact criteria and useful proofs}
\label{sec:squares}

\subsection{A complete criterion for direct denesting}
Let
\begin{equation}\label{eq:quadratic}
\alpha=\sqrt{a+b\sqrt r},\qquad a,b,r\in\Q,
\end{equation}
where $r>0$ is not a rational square, $b\ne0$, and $a+b\sqrt r>0$. Set
\begin{equation}
D=a^2-b^2r.
\end{equation}
The term $D$ is the norm of $a+b\sqrt r$ from $\Q(\sqrt r)$ to $\Q$. The following rational version of the classical criterion captures more than the familiar two-term ansatz \cite{bfht,gkioulekas}.

\begin{theorem}[Direct square-root denesting]\label{thm:direct}
Under the preceding hypotheses, the following are equivalent:
\begin{enumerate}[label=(\roman*),itemsep=1pt]
\item $\alpha$ belongs to a field generated over $\Q$ by finitely many square roots of positive rational numbers.
\item $\alpha$ has a representation as a signed sum of two square roots of nonnegative rational numbers.
\item $D$ is a square in $\Q$.
\end{enumerate}
When $d=\sqrt D\in\Q$, an explicit formula is
\begin{equation}\label{eq:direct}
\alpha=\sqrt{\frac{a+d}{2}}+\sgn(b)\sqrt{\frac{a-d}{2}}.
\end{equation}
\end{theorem}
\begin{proof}
Assume first that $D=d^2$ with $d\ge0$. Since $b\ne0$ and $r$ is nonsquare, $D\ne0$. The positivity of $a+b\sqrt r$, together with $a^2>b^2r$, forces $a>|b|\sqrt r>0$. Thus $x=(a+d)/2$ and $y=(a-d)/2$ are nonnegative, $x\ge y$, and
\[
x+y=a,\qquad xy=b^2r/4.
\]
Squaring $\sqrt x+\sgn(b)\sqrt y$ gives $a+b\sqrt r$, and this expression is nonnegative. This proves (iii)$\Rightarrow$(ii); (ii)$\Rightarrow$(i) is immediate.

For the converse, put $L=\Q(\alpha)$ and $F=\Q(\sqrt r)$. Since $\sqrt r=(\alpha^2-a)/b$, we have $F\subseteq L$, and $[L:\Q]$ is two or four. By (i), $L$ is a subfield of a multiquadratic extension. Such an extension is Galois with elementary abelian $2$-group; consequently $L/\Q$ is also Galois and all its automorphisms have order at most two.

If $L=F$, write $\alpha=u+v\sqrt r$ and let $\beta=u-v\sqrt r$. Then $\alpha\beta\in\Q$ and
\[
(\alpha\beta)^2=(a+b\sqrt r)(a-b\sqrt r)=D.
\]
If $[L:\Q]=4$, its Galois group is the Klein four-group. Let $\kappa$ be the nonidentity automorphism fixing $F$; then $\kappa(\alpha)=-\alpha$. Choose $\tau$ that changes the sign of $\sqrt r$ and set $\beta=\tau(\alpha)$. Since $\tau^2=1$, the product $\alpha\beta$ is fixed by $\tau$. Since $\kappa$ and $\tau$ commute, $\kappa(\beta)=-\beta$, so the product is also fixed by $\kappa$. It is therefore rational. Its square is again $D$.
\end{proof}

The necessity proof is worth retaining: unsuccessful coefficient matching in a chosen ansatz is not generally a proof of non-denestability. Here the field structure proves that no more elaborate multiquadratic expression can evade the two-term criterion.

If $b=0$, or if $r$ is a rational square, the expression is already a square root of a rational number and should be handled before applying the theorem. These elementary degeneracies are often omitted in informal formulas.

\subsection{Depth and degree are independent measures}
For $\alpha=\sqrt2+\sqrt3$,
\[
(\alpha^2-5)^2=24,
\qquad \alpha^4-10\alpha^2+1=0.
\]
Moreover $1/\alpha=\sqrt3-\sqrt2$, so
\[
\sqrt3=\frac{\alpha+\alpha^{-1}}2,
\qquad
\sqrt2=\frac{\alpha-\alpha^{-1}}2.
\]
Thus $\Q(\alpha)=\Q(\sqrt2,\sqrt3)$ and the displayed quartic is minimal. Its three quadratic subfields are $\Q(\sqrt2)$, $\Q(\sqrt3)$, and $\Q(\sqrt6)$. Denesting has changed the presentation of the field, not its degree.

This explicitly refutes the assertion, found in one supplied report, that direct denesting is equivalent to the result lying in a \emph{single quadratic extension} of the base field. The correct statement allows a multiquadratic field; even the standard example has degree four. Landau's expository discussion of this number gives a useful wider perspective \cite{landau98}.

\subsection{Fourth roots can do something square roots cannot}
The real degree-two nesting problem is not exhausted by Theorem~\ref{thm:direct} once fourth roots are admitted. Borodin et al. establish that, for the expression $\sqrt{a+b\sqrt r}$ over a real base field, arbitrary real unnested radicals give no additional possibilities beyond square roots and the relevant fourth root. In the rational setting above, the second test is
\begin{equation}\label{eq:indirecttest}
-rD=r(b^2r-a^2)\quad\hbox{is a square in }\Q.
\end{equation}
Thus real depth-one denesting exists precisely when the direct test or this indirect test succeeds \cite{bfht,gkioulekas}. The real-field hypothesis and the specified shape of the input are essential; this is not an unrestricted theorem about all nested expressions of every depth.

When $D<0$ and \eqref{eq:indirecttest} holds, define
\[
q=\sqrt{-D/r}\in\Q,\qquad
U=\frac{b+q}{2},\qquad V=\frac{b-q}{2}.
\]
The positivity of the original radicand forces $b>0$, while $q^2=b^2-a^2/r\le b^2$. Hence $U,V\ge0$, and an explicit identity is
\begin{equation}\label{eq:indirect}
\sqrt{a+b\sqrt r}=\sqrt[4]r\,\bigl(\sqrt U+\sgn(a)\sqrt V\bigr).
\end{equation}
Its square is $b\sqrt r+\sgn(a)\sqrt r\sqrt{b^2-q^2}=a+b\sqrt r$; its sign is nonnegative because $U\ge V$. This directly verifies the construction and its branch.

For example,
\begin{equation}
\sqrt{4+3\sqrt2}=\sqrt[4]2(1+\sqrt2)
=\sqrt[4]2+\sqrt[4]8.
\end{equation}
Here $D=-2$ is not a rational square, but $-rD=4$ is. Conversely,
\[
\sqrt{1+\sqrt2}
\]
has $D=-1$ and $-rD=2$; neither test succeeds. In the specified real depth-one model it cannot be denested, even though allowing an arbitrary collection of real root indices may initially look more powerful.

\subsection{Recursive identities and their limitations}
A typical practical step rewrites $\sqrt{X+Y}$ through $X^2-Y^2$:
\begin{equation}\label{eq:recursive}
\sqrt{X+Y}
=\sqrt{\frac{X+\sqrt{X^2-Y^2}}2}
+\sqrt{\frac{X-\sqrt{X^2-Y^2}}2},
\end{equation}
for real $X\ge Y\ge0$. Negative $Y$ requires a signed difference and a nonnegativity check. Merely applying this identity need not simplify anything: the new discriminant can be harder than the original radicand.

The practical issue is to find a useful split and accept it only when a well-founded complexity measure decreases. Jeffrey--Rich develops such a recursive method, and its influence is visible in modern source code \cite{jeffrey,sympysrc,raddenestsrc}. For example,
\begin{equation}\label{eq:depth3}
\sqrt{16-2\sqrt{29}+2\sqrt{55-10\sqrt{29}}}
=\sqrt5+\sqrt{11-2\sqrt{29}}.
\end{equation}
Squaring the right side proves the identity, and both terms are positive. The outer layer disappears, but the remaining square root is still nested.

Arithmetic-operation bounds in algorithms for prescribed quadratic towers do not directly bound the size of fully expanded expressions. Nor does the success of a top-level split establish global minimum depth. Borodin et al. explicitly distinguish tower-based procedures and their more general depth-two treatment \cite{bfht}. A production implementation should keep those guarantees separate.

\section{From radical bases to general denesting algorithms}
\label{sec:theory}

\subsection{Why linear independence is foundational}
Suppose an expression is expanded into products of radicals. Collecting like terms requires deciding which products are genuinely linearly independent over the base field. Formal monomials are not automatically a basis: $\sqrt2\sqrt8=4$, and adjoining a radical already present in the field contributes no new degree.

Besicovitch, Mordell, Siegel, and Kneser provide classical results on independence and degrees of radical extensions \cite{besicovitch,mordell,siegel,kneser}. Their hypotheses must be preserved. A formula such as
\[
[K(\sqrt[n_1]{a_1},\ldots,\sqrt[n_s]{a_s}):K]=n_1\cdots n_s
\]
requires independence assumptions; it is not a general rule. These results explain why many symbolic algorithms seek a reduced family of radical generators before simplification.

A relevant addition to the supplied bibliography is Perucca's 2025 preprint on linear relations among radicals. It describes the discrepancy between extension degree and the index of a multiplicative radical group, identifying the role of roots of unity and relations involving individual radicals. This is structural input for basis construction, not by itself a minimum-depth algorithm \cite{perucca}. Lenstra's \emph{Entangled radicals} lectures supply complementary background on interactions among radical extensions \cite{lenstra06}.

\subsection{Caviness--Fateman and the normalization tradition}
Fateman's thesis and Caviness--Fateman's SYMSAC paper belong to the history of exact symbolic simplification. The latter develops RADCAN-style treatment of radical expressions, especially unnested radical extensions \cite{fateman72,caviness76}. Its relevance is methodological: a linear basis and a normalization procedure can expose cancellation and make equality testing reliable.

It should not be retrospectively described as a general denesting algorithm. A canonical representative in a chosen field basis need not be the shortest or shallowest radical formula. This distinction remains visible in Maple's documented \texttt{radnormal} pipeline: basis construction, algebraic normalization, then conversion back to radicals \cite{maplerad}.

\subsection{Zippel's structural viewpoint}
Zippel's 1977 thesis and 1985 paper shift the subject from isolated identities toward algebraic dependence, basis construction, and a structural criterion for denesting \cite{zippel77,zippel85}. Rather than repeatedly guessing coefficients, the algorithmic task is to understand how an element sits inside a radical extension.

Landau's short 1992 note is an essential companion. It repairs a gap in the earlier argument and establishes necessity of the relevant Zippel condition, not just sufficiency \cite{landauNote}. Consequently a presentation that quotes a broad ``Zippel theorem'' without specifying its field hypotheses, its notion of denesting, and the later correction is incomplete. The 1985 paper remains a principal source, but should be read together with the note rather than through a compressed formula from a secondary survey.

\subsection{Landau's general framework}
Landau's FOCS 1989 paper, expanded in SIAM in 1992, is the central general denesting reference. Over a base field containing all roots of unity, it gives necessary and sufficient conditions and constructs a denesting when possible. Without that hypothesis, the stated result is within one of optimal after adjoining a single root of unity; the symbol $\zeta_\ell$ is used directly, not expanded into a nested radical. The algorithms require a splitting field of the minimal polynomial and have exponential running time in the general setting \cite{landau92}.

This is a substantial theoretical resolution, but not an unconditional promise of a cheap, shortest, branch-correct real expression over $\Q$ in every output language. Three choices remain visible: which cyclotomic constants are allowed, whether their construction contributes depth, and which radical branches are admitted. The expository article \emph{How to tangle with a nested radical} is the best preliminary reading before the technical paper \cite{landau94}.

For orientation, a radical tower can be written
\[
K=K_0\subseteq K_1\subseteq\cdots\subseteq K_d,
\qquad K_{j+1}=K_j(\theta_{j,1},\ldots,\theta_{j,m_j}),
\]
where powers of each newly adjoined element lie in $K_j$. A single layer can contain several parallel adjunctions. Replacing all of them by a chain of one-generator extensions artificially inflates the apparent depth. Galois closure and roots of unity make the multiplicative structure of these layers accessible to group theory, but constructing that closure can dominate the computation.

\subsection{Horng--Huang and minimum-depth radical solutions}
Horng--Huang treats both simplification of nested radicals and construction of radical solutions of solvable polynomials, with special attention to pure root extensions and minimum depth \cite{horng90,horng99}. The later paper specifically includes roots of unity in its model. The derived series of a solvable Galois group is central to the connection between successive radical layers and group structure.

One should not reduce this to the slogan ``minimum depth equals the derived length of the original Galois group over $\Q$'' without qualifications. Cyclotomic base changes, the definition of a pure radical layer, the chosen element rather than the whole splitting field, and the status of roots of unity all matter. Horng's thesis supplies a longer treatment \cite{horngThesis}. Cotner's dissertation is another pertinent thesis-level lead on nesting depths; its existence and title were checked, but its detailed theorem statements were not independently audited here \cite{cotner}.

Landau--Miller's polynomial-time result concerns deciding solvability by radicals for a polynomial in the stated representation \cite{landauMiller}. It is not a theorem that a formula of minimum nesting depth can always be printed in polynomial time in the length of an arbitrary nested input expression. Decision, construction, and output representation must be kept distinct.

\subsection{Bl\"omer: avoid unnecessary field expansion}
Bl\"omer's FOCS 1992 work addresses depth-two denesting through real or bounded-degree radicals without first constructing the full minimal-polynomial/splitting-field machinery. It supplies structure and size bounds and some applications beyond depth two. The dissertation places denesting alongside sums of radicals and exact algebraic representation \cite{bloemer92,bloemer93}.

His bounded-degree work, first presented at ESA 1997 and expanded in \emph{Algorithmica} in 2000, asks for an optimal or near-optimal expression using $D$th roots when the input uses $d$th roots and $D$ is a multiple of $d$. The running time is polynomial in the \emph{description size of the splitting field}, not automatically in the original radical string \cite{bloemer00}. This qualification is one of the most consequential corrections to the supplied summaries.

The 1991 paper on computing sums of radicals and the 1998 probabilistic zero-test are separate contributions \cite{bloemer91,bloemer98}. They should not be merged into a single vaguely dated Monte Carlo result. A decision procedure for equality of explicitly represented sums can have a very different complexity from an algorithm for finding their simplest radical presentation.

\subsection{Four sizes that complexity statements must identify}
Write $s$ for the size of the input expression or circuit, $N$ for the degree of an ambient algebraic extension, $S$ for a splitting-field description size, and $L$ for the size of a requested output formula. These parameters need not be polynomially related. For instance, $\Q(\sqrt{p_1},\ldots,\sqrt{p_m})$ has degree $2^m$ for distinct primes, while the list of radicals has only $m$ entries apart from coefficient bit lengths.

Likewise, a directed acyclic graph can share subexpressions that an expanded tree repeats many times. Arithmetic in an extension field is not a constant-cost bit operation when coordinates and coefficients grow. A claim of polynomially many field operations therefore does not establish polynomial Turing complexity, and any formula-producing algorithm requires at least enough time to write its output.

These observations do not negate the classical theorems. They tell the implementer which theorem actually applies. A useful complexity summary should always state the input encoding, the cost model, the root-index restrictions, and whether output is a shared circuit, a tower, or an expanded formula.

\section{Ramanujan identities and restricted higher-root problems}
\label{sec:ramanujan}

\subsection{Identities as evidence of field structure}
Ramanujan's formulas are more than a collection of attractive examples: they pose questions about units, norms, conjugates, and the interaction of pure and cyclotomic extensions. The collected papers and Berndt's editions of the notebooks are the historical sources; Berndt--Chan--Zhang explains the Indian context and the role of radicals and units \cite{ramanujan,berndtIV,berndtV,berndt97,berndt98}. An identity obtained by taking a power of an expression is easy to generate; a classification of when the reverse operation is possible is a substantially different achievement.

For example,
\begin{equation}\label{eq:ram-cube}
\sqrt[3]{\sqrt[3]2-1}
=\frac{1-\sqrt[3]2+\sqrt[3]4}{\sqrt[3]9}.
\end{equation}
An exact verification needs no numerical guessing. Put $t=\sqrt[3]2$. Reduction modulo $t^3-2$ gives
\[
(1-t+t^2)^3=9(t-1).
\]
The numerator is positive, since $1-t+t^2=(t-1/2)^2+3/4$, and the real cube is injective. Thus \eqref{eq:ram-cube} follows. Division by an unnested cube root does not introduce another layer under the convention of Section~\ref{sec:scope}.

Two further classical examples are
\begin{align}
\sqrt{\sqrt[3]5-\sqrt[3]4}
 &=\frac{\sqrt[3]2+\sqrt[3]{20}-\sqrt[3]{25}}3,
 \label{eq:ram-five}\\
\sqrt{\sqrt[3]{28}-3}
 &=\frac{\sqrt[3]{98}-\sqrt[3]{28}-1}3.
 \label{eq:ram-twentyeight}
\end{align}
These are not consequences of the rational quadratic discriminant test: their inner field has cubic rather than quadratic structure. Honsbeek's criterion provides the appropriate decision framework \cite{honsbeek99,honsbeek05,sury}.

Shanks's \emph{Incredible identities}, Spohn's subsequent letter, and Bumby's revisitation are useful historical companions \cite{shanks,spohn,bumby}. In this family of literature, an appealing printed identity should still be checked for normalization and branch conventions. In particular, the publication year of Bumby's paper is 1987; the much later online publication date is not the date of the mathematical work.

\subsection{A rational quartic criterion for a square root of two cube roots}
Let $\alpha,\beta\in\Q^\times$, suppose $\beta/\alpha$ is not a rational cube, and suppose $\sqrt[3]\alpha+\sqrt[3]\beta>0$. In Honsbeek's real-radical setting, the expression
\[
\sqrt{\sqrt[3]\alpha+\sqrt[3]\beta}
\]
denests precisely when the quartic
\begin{equation}\label{eq:honsbeek}
F(s)=s^4+4s^3+8\frac{\beta}{\alpha}s-4\frac{\beta}{\alpha}
\end{equation}
has a rational zero. This is a specialized theorem, not an algorithm for arbitrary depth-two mixtures of radical indices. Honsbeek's 1999 report and 2005 thesis are separate bibliographic objects; Sury's lecture notes provide an accessible account \cite{honsbeek99,honsbeek05,sury}.

The constructive part can be checked directly. If $s\in\Q$ satisfies \eqref{eq:honsbeek}, put $f=\beta-s^3\alpha$ and
\[
E=-\frac{s^2}{2}\sqrt[3]{\alpha^2}
   +s\sqrt[3]{\alpha\beta}+\sqrt[3]{\beta^2}.
\]
Writing $u=\sqrt[3]\alpha$ and $v=\sqrt[3]\beta$, direct expansion yields
\[
E^2-f(u+v)
=\left(\frac{s^4\alpha}{4}+s^3\alpha+2s\beta-\beta\right)u=0.
\]
The hypothesis on $\beta/\alpha$ implies $f\ne0$; the displayed identity and $u+v>0$ then imply $f>0$. Thus $|E|/\sqrt f$ is a real depth-one formula for the desired square root. This elementary calculation establishes sufficiency and checks the formula; the converse is the nontrivial field-theoretic part of the cited theorem.

For $(\alpha,\beta)=(5,-4)$, one may take $s=-2$, giving $f=36$ and \eqref{eq:ram-five}; for $(\alpha,\beta)=(28,-27)$, $s=-3$ gives $f=729$ and \eqref{eq:ram-twentyeight}. In each case the sign is chosen for the nonnegative square root. Exchanging $\alpha$ and $\beta$ cannot change denestability, because the original expression is unchanged. A nonsymmetric parametrization of certificates must not be confused with a nonsymmetric mathematical property.

\subsection{Cardan polynomials and an exact quadratic-field cubic test}
Osler's Cardan-polynomial treatment and the related work of Witu\l a--S\l ota explain a useful algebraic mechanism \cite{osler,witula}. Define
\[
C_0(X,P)=2,\quad C_1(X,P)=X,\quad
C_n(X,P)=XC_{n-1}(X,P)-PC_{n-2}(X,P).
\]
Induction proves
\begin{equation}\label{eq:cardan}
C_n(u+v,uv)=u^n+v^n.
\end{equation}
Indeed, multiplying $u^{n-1}+v^{n-1}$ by $u+v$ creates two cross terms, and subtracting $uv(u^{n-2}+v^{n-2})$ cancels exactly those terms. For $n=3$, this becomes $C_3(X,P)=X^3-3PX$.

The following formulation gives a complete and readily implementable test for \emph{membership in the original real quadratic field}. It is an equivalent rescaling of the criterion studied by Cavallo \cite{cavallo}; the proof is included to clarify what the test does and does not establish.

\begin{theorem}[A real cubic radical in a fixed quadratic field]\label{thm:cubic}
Let $a,b,p\in\Q$, with $b\ne0$ and $p>0$ nonsquare. Set
\[
N=a^2-b^2p.
\]
There exist $A,B\in\Q$ such that
\[
\sqrt[3]{a+b\sqrt p}=A+B\sqrt p
\]
if and only if $N=q^3$ for some $q\in\Q$ and the polynomial
\begin{equation}\label{eq:cubic-test}
T(X)=X^3-3qX-2a
\end{equation}
has a rational root $t$. When these conditions hold,
\begin{equation}\label{eq:cubic-solution}
A=\frac t2,\qquad B=\frac{b}{t^2-q}.
\end{equation}
\end{theorem}
\begin{proof}
Suppose first that $z=A+B\sqrt p$ has cube $a+b\sqrt p$, and put $z'=A-B\sqrt p$. Then $(z')^3=a-b\sqrt p$, so $q=zz'=A^2-B^2p$ is rational and satisfies $q^3=N$. Also $t=z+z'=2A$ is rational, and \eqref{eq:cardan} gives $t^3-3qt=2a$.

Conversely, assume $q^3=N$ and $T(t)=0$. The identity
\begin{align}
(t^2-q)^2(t^2-4q)
 &=t^2(t^2-3q)^2-4q^3\notag\\
 &=4a^2-4N=4b^2p
 \label{eq:cubic-certificate}
\end{align}
shows that $t^2-q\ne0$. With $A,B$ defined by \eqref{eq:cubic-solution}, division by $4(t^2-q)^2$ gives
\[
B^2p=A^2-q.
\]
Consequently
\begin{align*}
A^3+3AB^2p&=4A^3-3Aq=(t^3-3qt)/2=a,\\
3A^2B+B^3p&=B(4A^2-q)=B(t^2-q)=b.
\end{align*}
Thus $(A+B\sqrt p)^3=a+b\sqrt p$. Both sides define real quantities and the real cube is injective, proving the required branch equality.
\end{proof}

The discriminant of $T$ is $-108b^2p<0$, so $T$ has a unique real root. A rational root is therefore unique as well. Degenerate inputs with $b=0$ or square $p$ should be reduced separately rather than forced through the denominator in \eqref{eq:cubic-solution}.

For example, $(a,b,p)=(7,5,2)$ gives $q=-1$, $T(X)=X^3+3X-14$, and $t=2$. Hence
\[
\sqrt[3]{7+5\sqrt2}=1+\sqrt2.
\]
Similarly $\sqrt[3]{85+62\sqrt7}=1+2\sqrt7$. In contrast, for $a=2$, $b=1$, $p=3$, the norm is the cube $1$, but $X^3-3X-4$ has no rational root: the candidates $\pm1,\pm2,\pm4$ all fail. The cube-norm condition alone is insufficient. This proves nonmembership in $\Q(\sqrt3)$, not a prohibition on every other form of unnested radical expression.

\subsection{Why integer factorization is not the right complexity conclusion}
Cavallo's preprint contains an assertion connecting the cubic reducibility step with the integer-factorization problem, which one supplied report turns into a strong computational equivalence \cite{cavallo}. The mathematical criterion survives, but that complexity inference is not justified.

For the restricted problem of Theorem~\ref{thm:cubic}, the following is a deterministic polynomial-bit-time route. Write $N$ as a reduced rational number, and test whether its numerator and denominator are perfect cubes using exact integer-root computation. No prime factorization is needed. If the test succeeds, factor the degree-three rational polynomial $T$. Factoring polynomials over $\Q$ is polynomial-time by Lenstra--Lenstra--Lov\'asz; clear denominators and use that algorithm to detect a linear factor \cite{lll}. Finally reconstruct $A,B$ and verify the two rational coefficient equalities in the proof.

Enumerating divisors of an integer coefficient is one possible way to search for rational roots, but the existence of that implementation does not imply that integer factorization is necessary, or that rational-root testing is equivalent to it. The correct conclusion is that the \emph{specified quadratic-field membership problem lies in polynomial time}. Nothing here establishes a separation from integer factorization, nor a polynomial-time algorithm for all denesting problems.

\subsection{Pure cubic extensions, cyclic fields, and units}
Krepkii--Pimenov explains Ramanujan cube-root formulas through normal closures of pure cubic extensions over cyclic cubic fields. Their argument illustrates Hilbert's Theorem~90 and admits a prime-degree generalization. Antipov--Pimenov subsequently describes Ramanujan-type cubic formulas in a prescribed pure-cubic setting: the abstract identifies a bound of three summands and a cube-norm requirement, together with a qualified uniqueness statement \cite{krepkii,antipov}. Those statements should not be expanded into a universal classification of all cube-root expressions over arbitrary fields.

Shevelev's Ramanujan cubic polynomials and Barbero--Cerruti--Murru--Abrate's link with Shanks polynomials add a complementary source of identities through cyclic cubic fields, trigonometric expressions for roots, and Gaussian periods \cite{shevelev,barbero}. These works are particularly relevant when a radical identity is entangled with values of trigonometric functions at rational multiples of $\pi$, rather than presented initially as a small rational surd.

Dokmeci's MIT PRIMES report is a useful additional route into field-extension arguments and real radical denesting \cite{dokmeci}. It is retained as student research, not silently upgraded to a peer-reviewed general theorem. Readers should compare its conventions and hypotheses with the classical sources before treating apparently different criteria as equivalent.

\subsection{Finite constructions beyond cubes and results through 2026}
Campbell--Zujev's \emph{Variations on Ramanujan's nested radicals} concerns finite identities and constructions, including quartic and cubic examples; it should not be filed exclusively with infinite radicals \cite{campbell}. One representative identity is
\begin{equation}\label{eq:quartic-quotient}
\sqrt[4]{\frac{3+2\sqrt[4]5}{3-2\sqrt[4]5}}
=\frac{\sqrt[4]5+1}{\sqrt[4]5-1}
=\frac{3+\sqrt[4]5+\sqrt5+\sqrt[4]{125}}2.
\end{equation}
To verify the first equality, set $t=\sqrt[4]5$ and reduce
\[
(t+1)^4(3-2t)-(t-1)^4(3+2t)
\]
modulo $t^4-5$; the remainder is zero. Since $1<t<3/2$, all relevant denominators and the claimed fourth root are positive. The second equality follows by multiplying by $t-1$ and using $t^4=5$. Thus both algebraic equality and branch selection have been checked.

Girstmair's \emph{Reducing radicals in the spirit of Euclid}, published in 2021 after a 2020 preprint, explores representations of radicals of odd index $p\ge3$ associated with polynomials of degree $2p$ using irrationalities of degree at most $p$ \cite{girstmair}. The objective is degree reduction; it does not automatically coincide with reducing the number of nested root signs.

Osipov--Kytmanov's CASC 2021 chapter develops an algorithmic theory for triply nested real radicals over $\Q$ and supplies examples that cannot be simplified within its setting \cite{osipov21}. Its scope is more specific than a complete practical solver for unrestricted depth and arbitrary indices. Tonien's 2025 article constructs fifth-root identities through elliptic curves and includes computational material in SageMath \cite{tonien}; it is a source of higher-index families, not a blanket decision procedure.

A further addition is Osipov's 2026 paper on Ramanujan identities with nested cubic radicals. The publisher's abstract describes a universality result for specified cubic families over $\Q$ and limitations of analogous statements over arbitrary real base fields \cite{osipov26}. Because the accessible abstract renders some formulas poorly and the full text required authentication, this report records the verified bibliographic lead and scope rather than reconstructing an exact theorem from damaged notation.

\section{Equality testing, succinct representations, and complexity}
\label{sec:complexity}

\subsection{Equality is not order comparison}
An exact denester needs to verify identities, but an identity tester need not produce a denested expression. Likewise, testing whether an explicitly represented sum of square roots is zero differs from determining its sign. Bl\"omer's work gives polynomial-time control of the equality problem for explicitly encoded sums of radicals; the square-root-sum comparison problem has a different complexity history \cite{bloemer91,allender}.

A frequent misleading inference is that exact equality is generally undecidable because a simplifier sometimes returns an unevaluated expression. Finite algebraic expressions with specified algebraic branches can be represented and compared exactly using minimal polynomials, field arithmetic, and root isolation. The difficult questions are the cost of doing so from succinct input, and whether a particular restricted symbolic routine implements a complete decision procedure. General symbolic expressions involving transcendental functions constitute a different setting and should not be imported into this claim.

\subsection{What Radical Identity Testing actually takes as input}
Balaji--Nosan--Shirmohammadi--Worrell studies a polynomial circuit $f\in\Z[x_1,\ldots,x_k]$ evaluated at real radical constants:
\[
f\bigl(\sqrt[d_1]{a_1},\ldots,\sqrt[d_k]{a_k}\bigr)=0,
\]
where the nonnegative integers $a_i$ and positive root indices $d_i$ are binary encoded. The result places RIT in coNP under the Generalized Riemann Hypothesis. For the restriction to square roots of primes, it gives coRP under GRH and coNP unconditionally \cite{rit}.

The circuit has arithmetic gates applied to unnested radical constants. It is \emph{not} an arbitrary circuit with a root gate permitted at every internal node. The reported upper bounds therefore must not be transplanted to every conceivable nested-radical language. Nor does membership in a decision complexity class yield a minimum-depth formula constructor. The LICS 2022 publication and the revised 2024 arXiv version belong to one research line, not two independent results.

\subsection{Power-series methods and square-root sums}
Gaillard--Jindal relates orders of linear combinations of power series to Wronskians, sharpens the understanding of the Wronskian method, and examines a succinct-input generalization of the equality variant of square-root sums. Under GRH, its second part reduces that generalization to a one-dimensional variant \cite{gaillard}. This is not a resolution of the unrestricted sign-comparison problem.

Allender et al. provides important background on numerical decision problems and counting-hierarchy bounds \cite{allender}. It is included as a foundational complexity reference, not labelled an automatically up-to-date ``best known bound'' without a separate exhaustive frontier review. For a denesting implementation, these papers are most useful when choosing certificate formats and input representations, rather than as ready-made rewriting algorithms.

\section{Computer-algebra implementations and their guarantees}
\label{sec:cas}

\subsection{A capability map, not a benchmark ranking}
The systems below expose different abstractions. Some transform radical syntax; others normalize an algebraic number; others supply the field machinery from which a denester could be built. The entries describe documented behavior as checked for this report. Only the SymPy examples in Section~\ref{sec:tests} were executed here. Examples for other systems are illustrative usage, not measured results.

\begin{longtable}{@{}>{\raggedright\arraybackslash}p{0.19\textwidth}>{\raggedright\arraybackslash}p{0.34\textwidth}>{\raggedright\arraybackslash}p{0.40\textwidth}@{}}
\toprule
System & Relevant interface & What should not be inferred \\
\midrule\endfirsthead
\toprule System & Relevant interface & What should not be inferred \\
\midrule\endhead
SymPy & \texttt{sqrtdenest}; \texttt{nthroot}; exact algebraic-number tools & \texttt{max\_iter=3} is not a depth-three theorem; failure is not a general impossibility certificate.\\[4pt]
Maxima & \texttt{raddenest}; \texttt{radcan}; legacy square-root helpers & The port has higher-root extensions, but is not documented as a complete general minimum-depth implementation.\\[4pt]
Wolfram Language & \texttt{RootReduce}, \texttt{ToRadicals}; repository \texttt{RadicalDenest} & A canonical \texttt{Root} object is not an unnested radical; the repository function is heuristic.\\[4pt]
Maple & \texttt{radnormal}, \texttt{simplify(...,radical)}, \texttt{evala} & Exact radical normalization and zero detection do not imply minimum radical depth.\\[4pt]
SageMath & \texttt{AA}/\texttt{QQbar}; \texttt{radical\_expression()} & The conversion can return the algebraic number unchanged; radical expansion need not improve depth.\\[4pt]
PARI/GP, Magma & Field factorization, subfields, embeddings, optimized field presentations & A smaller defining polynomial does not by itself give a shallower expression for a specified element.\\[4pt]
FriCAS/Axiom & \texttt{RadicalSolvePackage} and simplification facilities & Solving a low-degree polynomial by radicals is not a proof of optimal denesting.\\
\bottomrule
\end{longtable}
Sources for the rows are the official documentation and source-code records in the software bibliography \cite{sympydocs,sympysrc,raddenestsrc,wolframroot,wolframradical,wolframrd,maplerad,maplesimp,sageaa,sagesymbolic,pari,magma,fricas}.

\subsection{SymPy: specialized square roots, plus a real nth-root routine}
The public \texttt{sqrtdenest(expr,max\_iter=3)} entry point repeatedly applies specialized transformations. Its documentation explicitly cites Borodin et al. and Jeffrey--Rich. The source separates matching and recursive work, including helpers such as \texttt{\_sqrtdenest0}, \texttt{\_sqrt\_match}, and \texttt{\_denester}; their private names are useful navigation aids, not stable public APIs \cite{sympydocs,sympysrc}.

The iteration parameter bounds repeated passes. It neither limits the input to depth three nor promises a globally minimum result at any permitted depth. A heuristic can act on a deep expression when local patterns are favorable and leave a shallow one unchanged when its supported methods do not apply. Historical issues are useful regression-test leads, but an old issue is not evidence of a bug in the version executed here \cite{sympyissue}.

The supplied reports mostly overlooked \texttt{nthroot}, available from the simplification module. Its import and signature are:
\begin{lstlisting}
from sympy.simplify.simplify import nthroot
nthroot(expr, n, max_len=4, prec=15)
\end{lstlisting}
 The documented routine seeks the \emph{real} $n$th root of a sum of surds, first using candidate recognition and then a minimal-polynomial route when needed \cite{sympydocs}. It succeeds on the cube- and fifth-root examples below. It is therefore inaccurate to infer that the whole library cannot denest cube roots merely because \texttt{sqrtdenest} is specialized for square roots.

Keep three other operations separate. \texttt{radsimp} primarily handles radical denominators; \texttt{powdenest} simplifies powers under its assumptions; algebraic-number normalization provides exact equality and field representation. Aggressive forced power transformations may change branch semantics, so they cannot serve as a general proof of a denesting identity \cite{sympydocs}.

\subsection{Maxima: the port and its extensions}
Gilles Schintgen's \texttt{raddenest} source identifies its ancestry in SymPy 1.0 but also contains original extensions. Its Cardan-style helper treats certain $n$th roots in quadratic fields, and its Ramanujan helper uses the rational-quartic approach for a square root of a sum of cube roots. The source and tests are more informative than describing the package simply as a port \cite{raddenestsrc}.

A typical session begins with \texttt{load("raddenest")}; real-domain settings and the treatment of powers of negative numbers should be made explicit. The package attempts recursive simplification rather than building a full splitting field for every expression. \texttt{radcan} instead belongs to the radical-normalization tradition. The old \texttt{sqdnst}/\texttt{sqrtdenest} names remain relevant when reading historical Maxima examples, but version-specific availability should be checked in the installed manual rather than assumed from an old mailing-list post \cite{maximamanual,maximathread}.

The inspected higher-root code sometimes enumerates rational-root candidates through integer divisors. That is an implementation choice, not evidence that the underlying restricted decision problem is integer-factorization-hard. Nor do the examples justify an unmeasured claim that this package is the ``most capable'' free denester in every input class.

\subsection{Wolfram Language: exact algebraic objects versus heuristic denesting}
\texttt{RootReduce} transforms exact algebraic combinations into canonical algebraic-number representations. It can expose rational or low-degree simplifications, but the canonical object may be \texttt{Root}, not a radical expression \cite{wolframroot}. \texttt{ToRadicals} converts low-degree algebraic roots to radicals: the documentation guarantees degree at most four, while some higher-degree cases may fail even when a radical expression exists. Parameterized transformations have additional qualification about exceptional parameter values \cite{wolframradical}.

The separately contributed \texttt{ResourceFunction["RadicalDenest"]} explicitly uses heuristics, supports a time constraint, and returns the best result found. The page inspected lists a five-second default and version 2.1.0 dated 29 April 2024. Its author is Swastik Banerjee, with work by Corey Ziegler, Bill Gosper, and Daniel Lichtblau acknowledged \cite{wolframrd}. A resource function and a built-in kernel primitive should not be presented as the same product feature.

\begin{lstlisting}
RootReduce[Sqrt[5 + 2 Sqrt[6]]]
ToRadicals[RootReduce[expr]]
ResourceFunction["RadicalDenest"][expr, TimeConstraint -> 10]
\end{lstlisting}
These are usage patterns, not a guarantee that the second pipeline minimizes radical depth. \texttt{FullSimplify} may choose a different expression according to its own complexity objective. \texttt{PowerExpand} should not be used to bypass the branch check needed for a proof. The associated Wolfram Community discussions are useful accounts of heuristic development, not substitutes for theorems about completeness \cite{wolframcommunity}.

\subsection{Maple: what \texttt{radnormal} actually promises}
Maple documents \texttt{radnormal} as normalization for radical numbers. For algebraic-number expressions in its stated domain it detects zero exactly, while denesting occurs only in some cases. The documented pipeline constructs a basis using indexed \texttt{RootOf} representations, applies algebraic normalization, and converts back to radicals. It uses principal roots \cite{maplerad}.

This is a stronger and more useful description than guessing at proprietary internals. It does not establish that Maple implements the entire Landau or Bl\"omer algorithm, nor that it returns minimum-depth expressions. \texttt{simplify(...,radical)}, square-root simplification, radical combination, and \texttt{evala} have their own roles and assumptions \cite{maplesimp}. The single command name ``simplify'' cannot be used to erase distinctions among normalization, rationalization, solving, and denesting.

\subsection{SageMath, PARI/GP, Magma, and FriCAS}
Sage's real and complex algebraic fields, \texttt{AA} and \texttt{QQbar}, attach exact algebraic information and a chosen embedding to a number. The \texttt{radical\_expression()} method attempts conversion to a symbolic radical expression and may return the original algebraic number when unsuccessful. Returning that exact object is not a numerical approximation and not a non-denestability certificate \cite{sageaa}. Symbolic \texttt{canonicalize\_radical()} has a separate normalization role, with documentation warning about assumptions in radical simplification \cite{sagesymbolic}.

PARI/GP exposes \texttt{nfsubfields}, \texttt{nfisincl}, \texttt{nffactor}, and related field tools \cite{pari}. For the standard example, a useful exploration is
\begin{lstlisting}
f = x^4 - 10*x^2 + 1;
K = nfinit(f);
nfsubfields(K, 2)
\end{lstlisting}
The mathematical answer consists of the three quadratic subfields described earlier. An implementation still has to reconstruct the chosen element in suitable generators and certify its real branch. Field discovery is therefore only part of a denester.

Magma's optimized number-field representations aim at better defining data, not automatically a minimal-depth display of every element \cite{magma}. FriCAS/Axiom's \texttt{RadicalSolvePackage} exposes radical polynomial solving, which is adjacent to rather than synonymous with denesting \cite{fricas}. Giac/Xcas and the newer FriCAS simplification names in the supplied reports are retained as investigation leads; this survey does not turn incomplete documentation searches into categorical claims that these systems lack all denesting capability. Honda's Maxima \texttt{GaloisGroupSolver} is another adjacent source for solvable-polynomial constructions \cite{honda}.

\section{Reproducible examples and implementation lessons}
\label{sec:tests}

\subsection{What was executed}
The accompanying \texttt{verify\_examples.py} was run with SymPy 1.14.0. It checks the following transformations; equivalent output factorizations are normalized here for readability.
\begin{longtable}{@{}>{\raggedright\arraybackslash}p{0.47\textwidth}>{\raggedright\arraybackslash}p{0.46\textwidth}@{}}
\toprule Input & Result \\
\midrule\endfirsthead
\toprule Input & Result \\
\midrule\endhead
$\sqrt{5+2\sqrt6}$ & $\sqrt2+\sqrt3$\\[5pt]
$\sqrt{5-2\sqrt6}$ & $\sqrt3-\sqrt2$\\[5pt]
$\sqrt{4+3\sqrt2}$ & $\sqrt[4]2(1+\sqrt2)$\\[5pt]
$\sqrt{1+\sqrt2}$ & Unchanged by \texttt{sqrtdenest}\\[5pt]
Expression on the left of \eqref{eq:depth3} & $\sqrt5+\sqrt{11-2\sqrt{29}}$\\[5pt]
\texttt{nthroot} of $85+62\sqrt7$, $n=3$ & $1+2\sqrt7$\\[5pt]
\texttt{nthroot} of $41-29\sqrt2$, $n=5$ & $1-\sqrt2$\\
\bottomrule
\end{longtable}
The fourth row has an independent impossibility proof in the real depth-one model from Section~\ref{sec:squares}; that conclusion does not follow merely from the unchanged software output.

The script also checks the minimal polynomial of $\sqrt2+\sqrt3$, two polynomial-remainder certificates for Ramanujan-type identities, and Theorem~\ref{thm:cubic}. For the latter, it constructs all cases with
\[
p\in\{2,3,5,6,7\},\qquad -3\le A\le3,\qquad
B\in\{-3,-2,-1,1,2,3\},
\]
forms $a=A^3+3AB^2p$ and $b=3A^2B+B^3p$, and recovers $(A,B)$. All $5\cdot7\cdot6=210$ cases passed. There are additional positive and negative hand-picked tests. These computations corroborate an implementation of a proved criterion; they are not a substitute for the proof and not a performance comparison with other CAS products.

\subsection{An exact verification pattern}
A candidate denesting deserves two certificates. The first is algebraic equality: reduction in a common field or quotient ring, or an equivalent exact algebraic-number comparison. The second is selection of the correct root or embedding. For a positive square root, equality of squares is insufficient unless the candidate is nonnegative. For a real odd root, injectivity of the real power gives the branch conclusion directly. For complex radicals, root-of-unity factors must be tracked.

Numerical recognition is useful for generating candidates but not certifying them. Equal decimal approximations can result from cancellation or inadequate precision. Likewise, two expressions satisfying the same polynomial need not be equal: $\sqrt2$ and $-\sqrt2$ provide the simplest counterexample. A minimal polynomial must be supplemented by an isolating interval, isolating complex region, or an equivalent exact embedding specification.

For example, to prove \eqref{eq:fifth-real}, expand $(1-\sqrt2)^5$ in the basis $1,\sqrt2$, obtaining $41-29\sqrt2$. Since the candidate is real, it is the unique real fifth root. A transformation of the principal complex power would require a different branch factor; the same polynomial identity does not certify that transformation.

\subsection{A practical architecture suggested by the literature}
The following is an engineering synthesis, not a claimed new complete algorithm. Represent the input both as a shared expression graph and as an exact algebraic object with a selected embedding. Run cheap transformations first: rational-root extraction, square-root norm tests, supported Cardan cases, and the specialized Honsbeek quartic. Then attempt recursive square-root transformations and a reduced radical basis. Only when those fail should expensive subfield, Kummer, or splitting-field work be considered \cite{bfht,jeffrey,bloemer93,landau92}.

Keep the optimization objective explicit. One reasonable lexicographic cost is
\[
(\text{depth},\ \text{allowed-index penalty},\ \text{graph size},\
\text{coefficient bit size}).
\]
A different application may prioritize numerical stability over display depth. Denesting can create long sums with severe cancellation, and an exact \texttt{Root} representation may be better for evaluation. A user interface should distinguish ``no improvement found within the search budget'' from ``no expression exists in the specified model.'' Only the latter requires an impossibility certificate.

It is also useful to preserve an audit trail: original assumptions, selected branches, the transformation applied, the exact equality certificate, and the complexity comparison that caused the new form to be accepted. This prevents simplification cycles and makes a hybrid symbolic/algebraic implementation debuggable.

\section{Infinite nested radicals and other adjacent literature}
\label{sec:infinite}

\subsection{Convergence must precede evaluation}
For nonnegative $a_n$, consider finite truncations
\[
r_n=\sqrt{a_1+\sqrt{a_2+\cdots+\sqrt{a_n}}}.
\]
Herschfeld's theorem characterizes convergence by boundedness of the sequence $a_n^{2^{-n}}$ \cite{herschfeld}. It is not necessary that the radicands $a_n$ themselves be bounded. For example, $a_n=2^{2^n}$ gives $a_n^{2^{-n}}=2$ at every index. Borwein--de Barra supplies a later discussion of nested radicals and related analytic behavior \cite{borwein}.

Solving a fixed-point equation alone does not establish convergence of a proposed infinite nest or select its limit. For the simple positive nest $x=\sqrt{c+x}$ with $c>0$, one can separately show that zero-tail truncations increase and remain below the positive solution of $x^2-x-c=0$; only then does taking limits justify the fixed-point equation. Varying radicands or minus signs require their own analysis.

Mitra's generalizations to $n$th roots, Rahal's introductory thesis, and elementary expositions of Ramanujan's famous infinite radical are worthwhile side reading \cite{mitra,rahal,infiniteblogs}. They are indexed separately because a convergence theorem does not decide whether a fixed finite algebraic expression has a shallower radical representation.

\subsection{Algebraic equations arising from periodic nests}
Michele Elia studies algebraic equations associated with Ramanujan's work, including equations generated by iterated square-root patterns \cite{elia}. Eliminating a periodic nest can create a high-degree polynomial; one then has to determine which roots satisfy the original equations and their signs. This connects nested radicals with Galois theory, but the input is an implicit equation rather than a given finite radical expression to simplify.

Vi\`ete-type products, trigonometric nested radicals, and the classical literature surrounding $\pi$ provide additional context. Borwein--Borwein's \emph{Pi and the AGM} is a background source, not a general denesting manual \cite{piagm}. Similarly, Euclid's Book~X is historically relevant to constructions and quadratic irrationalities, while modern degree-reduction statements should be attributed to their modern proofs, such as Girstmair's \cite{euclid,girstmair}.

\section{Reading routes and research directions}
\label{sec:routes}

\subsection{A route from elementary algebra to complete theory}
For a first rigorous pass, read the proof in Section~\ref{sec:squares}, then Gkioulekas and Borodin et al. Follow with Landau's expository article, Zippel's 1985 paper together with Landau's correction, and finally Landau's technical paper. Horng--Huang is the next stop when the objective is minimum depth for radical solutions rather than one-layer simplification. Bl\"omer's thesis and bounded-degree paper provide the representation-sensitive algorithmic viewpoint \cite{gkioulekas,bfht,landau94,zippel85,landauNote,landau92,horng99,bloemer93,bloemer00}.

For implementation, Jeffrey--Rich should be read beside actual SymPy and Maxima source, not in isolation. Geddes--Czapor--Labahn, Zippel's textbook, Davenport--Siret--Tournier, and von zur Gathen--Gerhard supply the polynomial and symbolic-algebra background. Cohen gives the number-field tools; Cox and Newman supply radical extensions and Galois theory \cite{jeffrey,sympysrc,raddenestsrc,geddes,zippelbook,davenport,gathen,cohen,cox,newman}.

For Ramanujan-type work, begin with Berndt--Chan--Zhang and Sury, continue through Honsbeek and Osler, then compare the pure-cubic classifications with Cavallo's quadratic-field membership problem. Those are different families. Tonien supplies a route toward explicit fifth-root identities, while the 2021 and 2026 Osipov papers extend the recent specialized literature \cite{berndt97,berndt98,sury,honsbeek05,osler,antipov,cavallo,tonien,osipov21,osipov26}.

\subsection{Useful public discussions and search strategy}
Wikipedia and MathWorld are useful orientation and reference-discovery tools, but important scope restrictions should be checked against the original papers. Stan Brown's elementary page is useful for worked square-root examples. MathOverflow's discussion of the relevance of Landau's algorithm and Math StackExchange discussions of the discriminant criterion and Ramanujan identities help expose common conceptual confusions \cite{wiki,mathworld,brown,mo,mse}. These informal sources are not used here to replace primary evidence for the core theorems.

The most productive search vocabulary includes \emph{denesting}, \emph{simplification of radicals}, \emph{pure root extensions}, \emph{minimum radical depth}, \emph{radical independence}, \emph{Ramanujan-type formulas}, and \emph{radical identity testing}. Searching only for ``nested radical'' retrieves a large amount of infinite-radical material. Citation chaining from BFHT, Zippel, Landau, Bl\"omer, and Honsbeek is more selective. For recent work, distinguish an arXiv posting, a revised version, an online publication date, and the journal's issue year.

\subsection{Questions worth pursuing without overstating what is open}
The existence of broad theoretical denesting algorithms should not be hidden behind a claim that ``no general algorithm exists.'' The practical gap is between such algorithms and convenient, efficient, branch-correct implementations with explicit guarantees. In particular, it would be valuable to compare certificate-producing implementations on a shared exact corpus, investigate output-sensitive complexity with shared expression graphs, and expose restricted-model non-denestability certificates alongside heuristic search results.

Further directions include a carefully specified cost model for roots of unity; extending useful special-family methods without incurring a full splitting-field expansion; and connecting succinct equality certificates with actual expression construction. These are proposed research and engineering directions, not claims that each formulation is an independently verified open problem. A credible benchmark must include failures, mixed indices, negative real odd roots, parameter assumptions, coefficient growth, and the distinction between a compact field representation and a compact radical formula.

\appendix
\section{Consolidation audit and substantive corrections}
\label{sec:audit}

\subsection{Coverage of the supplied archive}
The archive contained reports labelled ChatGPT, Claude, Gemini, and Grok. Parallel Markdown, LaTeX, and PDF renderings were treated as versions of one report, not independent sources. Gemini's encoded equation images were reconciled with its LaTeX.

The contributions were complementary: ChatGPT supplied an extensive algorithmic backbone and field-tool discussion; Claude supplied specialized papers, theses, and software leads; Grok supplied historical organization and background references; Gemini supplied a cubic-field discussion and pedagogical leads. Material was merged by topic and bibliographic identity. Distinct papers were kept distinct; conference/journal versions and access mirrors were grouped. Copied encyclopedia mirrors, search-result scraps, package-listing pages, and unsupported implementation speculation were not elevated into independent mathematical sources.

\subsection{Corrections affecting the mathematical meaning}
\begin{description}[style=nextline,leftmargin=0pt,labelwidth=0pt]
\item[Quadratic versus multiquadratic.]
Direct denesting does not force the result into a single quadratic extension. The standard $\sqrt2+\sqrt3$ example has degree four. Theorem~\ref{thm:direct} supplies the correct field statement.
\item[Allowed root indices.]
A theorem about $\sqrt{a+b\sqrt r}$ over a real base field is not a theorem about all shapes and depths of nested square roots. Square-root-only and square-/fourth-root denesting are different output models.
\item[Cubic reducibility versus integer factorization.]
The supplied complexity equivalence is unsupported. Rational polynomial factorization gives a polynomial-bit-time route for the restricted problem in Theorem~\ref{thm:cubic}. Divisor enumeration is not a hardness reduction.
\item[An impossible asymmetry.]
The expressions $\sqrt{\sqrt[3]{-4}+\sqrt[3]5}$ and $\sqrt{\sqrt[3]5+\sqrt[3]{-4}}$ are identical. A report's assertion that one denests and the other does not is false; it confuses a parametrization with the underlying property.
\item[Field-size bounds.]
Polynomial dependence on splitting-field description size is not automatically polynomial dependence on the original radical expression. Minimum depth, minimum formula size, and minimum algebraic degree are not interchangeable.
\item[Identity testing.]
RIT uses arithmetic circuits evaluated at specified radical constants, not arbitrary internal radical gates. Equality, sign comparison, and finding a minimum-depth expression are separate tasks.
\item[Infinite radicals.]
Herschfeld's condition bounds $a_n^{2^{-n}}$, not $a_n$. Infinite-radical convergence and finite denesting require different arguments. Campbell--Zujev's finite constructions remain in the finite-radical section.
\end{description}

\subsection{Corrections affecting software claims}
SymPy's \texttt{max\_iter} is a pass limit, not a maximum nesting depth. Its \texttt{nthroot} routine invalidates a library-wide claim that only square roots can be denested. Maxima's \texttt{raddenest} is a port with substantive higher-root extensions, not just a renamed square-root routine. Conversely, no comparison performed here supports a universal superiority ranking among the systems.

Exact normalization in Maple, \texttt{RootReduce} in Wolfram Language, a Sage algebraic-number object, and optimized field presentations in Magma are not all ``denesting'' in the same sense. The absence of a named interface in a documentation search does not prove the absence of every implementation. Statements about proprietary algorithms, the supposed computational limitations of implementation languages, or current unfixed bugs were removed unless supported by inspected documentation or a versioned test.

\subsection{Bibliographic repairs and unresolved leads}
Landau's \emph{A note on ``Zippel denesting''} occupies pages 41--45, not 41--46. Horng--Huang's 1999 paper occupies pages 881--885, not 881--886. Bumby's paper is from 1987, notwithstanding its 2024 online publication date. Elia's paper is by Michele Elia, not Berndt and Andrews. Dokmeci's given name is Kaan. Honsbeek's 1999 report must not be cited with the title and date of the 2005 dissertation. The report also distinguishes Girstmair's 2020 preprint from the 2021 journal publication and Krepkii--Pimenov's 2015 issue date from its 2016 online date.

Some leads remain explicitly unverified at the capability or theorem level: the exact release history of legacy Maxima \texttt{sqdnst}; the reported newer FriCAS \texttt{rsimp}/\texttt{recursiveRootSimplification} interfaces; an isolated Rust \texttt{denest\_sqrt} mention; and assertions about Giac/Xcas's complete denesting coverage. They are not used to support the main conclusions. Unresolved generic question titles without stable identifiers were replaced, where possible, by concrete public discussion links. Informal infinite-radical blogs are retained only as ancillary exposition.

\section{Reproduction notes}
The ZIP contains this self-contained LaTeX source, its compiled PDF, and \texttt{verify\_examples.py}. The bibliography is embedded; no external \texttt{.bib} file, image download, or network access is needed to compile it. Use \texttt{latexmk -pdf radical\_denesting\_unified.tex}, or run pdfLaTeX twice.

The verification script requires Python 3.10 or later and SymPy; the test run used SymPy 1.14.0. Run \texttt{python verify\_examples.py}. Its exact cubic criterion and the meaning of a \texttt{None} result are documented in Section~\ref{sec:tests} and the source docstrings.

\clearpage
\phantomsection
\addcontentsline{toc}{section}{Annotated bibliography and source directory}
\renewcommand{\refname}{Annotated bibliography and source directory}
\begin{thebibliography}{99}
\setlength{\itemsep}{3pt}
\small
\bibgroup{I. Independence, normalization, and general algorithms}

\bibitem{besicovitch} A.~S. Besicovitch, ``On the linear independence of fractional powers of integers,'' \emph{Journal of the London Mathematical Society} \textbf{15} (1940), 3--6. \status{B}
Classical independence theorem underlying the use of radical monomials as field-basis elements. Retained as foundational background; its hypotheses should be read in the original rather than inferred from the title.

\bibitem{mordell} L.~J. Mordell, ``On the linear independence of algebraic numbers,'' \emph{Pacific Journal of Mathematics} \textbf{3} (1953), 625--630. \status{B}
A further classical source for independence of algebraic radicals. Relevant to exact basis construction, not a stand-alone denesting implementation.

\bibitem{siegel} C.~L. Siegel, ``Algebraische Abh\"angigkeit von Wurzeln,'' \emph{Acta Arithmetica} \textbf{21} (1972), 59--64. \doi{10.4064/aa-21-1-59-64}. \status{A}
Algebraic dependence of roots; part of the structural foundation used in later denesting algorithms.

\bibitem{kneser} M. Kneser, ``Lineare Abh\"angigkeit von Wurzeln,'' \emph{Acta Arithmetica} \textbf{26}(3) (1975), 307--308. \doi{10.4064/aa-26-3-307-308}. \status{A}
A short, important source on linear dependence and radical extensions. The issue year is recorded here as 1975.

\bibitem{fateman72} R.~J. Fateman, \emph{Essays in Algebraic Simplification}, PhD thesis, Massachusetts Institute of Technology, 1972. \status{B}
Historical background for algebraic simplification and the MACSYMA tradition. Distinguished from the later joint SYMSAC paper.

\bibitem{caviness76} B.~F. Caviness and R.~J. Fateman, ``Simplification of radical expressions,'' in \emph{Proceedings of the 1976 ACM Symposium on Symbolic and Algebraic Computation}, 329--338.
\access{https://people.eecs.berkeley.edu/~fateman/papers/radcan.pdf}{Author-hosted paper}. \status{T}
Foundational RADCAN-style normalization and exact treatment of radical expressions. Essential for understanding why canonicalization in a field basis and minimum-depth denesting are different tasks.

\bibitem{zippel77} R. Zippel, \emph{Simplification of Radicals with Applications to Solving Polynomial Equations}, master's thesis, Massachusetts Institute of Technology, 1977; related contribution, ``Radical simplification made easy,'' \emph{MACSYMA Users' Conference}, 1977. \status{B}
Early work preceding the 1985 journal article. These related historical records are grouped, without assuming that their texts are identical.

\bibitem{lll} A.~K. Lenstra, H.~W. Lenstra, Jr., and L. Lov\'asz, ``Factoring polynomials with rational coefficients,'' \emph{Mathematische Annalen} \textbf{261} (1982), 515--534. \doi{10.1007/BF01457454}. \status{T}
The polynomial-time rational polynomial-factorization result used to correct the alleged equivalence between the restricted cubic denesting test and integer factorization.

\bibitem{bfht} A. Borodin, R. Fagin, J.~E. Hopcroft, and M. Tompa, ``Decreasing the nesting depth of expressions involving square roots,'' \emph{Journal of Symbolic Computation} \textbf{1}(2) (1985), 169--188. \doi{10.1016/S0747-7171(85)80013-4}.
\access{https://www.cs.toronto.edu/~bor/Papers/decreasing-nesting-depth-for-expressions.pdf}{Author-hosted PDF}. \status{T}
Principal reference for square-root denesting, the role of fourth roots, real-field restrictions, and tower-based algorithms. Read the precise input shape and representation assumptions with each complexity statement.

\bibitem{zippel85} R. Zippel, ``Simplification of expressions involving radicals,'' \emph{Journal of Symbolic Computation} \textbf{1}(2) (1985), 189--210. \doi{10.1016/S0747-7171(85)80014-6}. \status{A}
Structural treatment of radical expressions and algebraic dependence. Read together with Landau's correction, not as an unqualified universal criterion detached from its field hypotheses.

\bibitem{landauMiller} S. Landau and G.~L. Miller, ``Solvability by radicals is in polynomial time,'' \emph{Journal of Computer and System Sciences} \textbf{30}(2) (1985), 179--208. \status{A}
Algorithmic Galois-theory infrastructure. Deciding radical solvability of a polynomial is not identical to constructing a shortest or shallowest formula for an already given radical expression.

\bibitem{landau92} S. Landau, ``Simplification of nested radicals,'' \emph{SIAM Journal on Computing} \textbf{21}(1) (1992), 85--110. \doi{10.1137/0221009}. Conference version: \emph{FOCS 1989}, 314--319, \doi{10.1109/SFCS.1989.63496}. \status{A}
The central general denesting paper. The journal abstract specifies the roots-of-unity hypothesis, the near-optimal statement when it is absent, the splitting-field requirement, and exponential general running time. Related versions are grouped here.

\bibitem{landauNote} S. Landau, ``A note on `Zippel denesting','' \emph{Journal of Symbolic Computation} \textbf{13}(1) (1992), 41--45. \doi{10.1016/0747-7171(92)90004-N}.
\access{https://web.cs.umass.edu/publication/docs/1990/UM-CS-1990-120.pdf}{UMass technical-report version}. \status{T}
Repairs the gap in the earlier argument and establishes necessity of the relevant Zippel condition. Page range corrected from a discrepancy in the supplied reports.

\bibitem{horng90} G. Horng and M.-D.~A. Huang, ``Simplifying nested radicals and solving polynomials by radicals in minimum depth,'' \emph{FOCS 1990}, beginning at p.~847.
\access{https://ieeexplore.ieee.org/document/89607/}{IEEE record, document 89607}. \status{A}
Also cited under the longer title ``On simplifying nested radicals and solving polynomials by pure nested radicals of minimum depth.'' Secondary and author bibliographies disagree between terminal pages 854 and 856; the identifier is more reliable than silently selecting one. Foundational minimum-depth work.

\bibitem{horngThesis} G. Horng, \emph{Simplifying Nested Radicals and Solving Polynomials by Nested Radicals in Minimum Depth}, PhD thesis, University of Southern California, 1990. \doi{10.5555/142931}. \status{A}
Longer thesis treatment associated with the Horng--Huang program. Useful for definitions of pure root extensions and the relationship with Galois-group structure.

\bibitem{bloemer91} J. Bl\"omer, ``Computing sums of radicals in polynomial time,'' \emph{FOCS 1991}, 670--677. \doi{10.1109/SFCS.1991.185434}. \status{A}
Exact computation and equality of explicitly represented sums of radicals. Kept distinct from the author's denesting and probabilistic zero-testing papers.

\bibitem{bloemer92} J. Bl\"omer, ``How to denest Ramanujan's nested radicals,'' \emph{FOCS 1992}, 447--456. \doi{10.1109/SFCS.1992.267807}. \status{A}
A key depth-two result avoiding unnecessary construction of full minimal-polynomial and splitting-field data. Relates structure and size bounds to denesting by real or bounded-degree radicals.

\bibitem{bloemer93} J. Bl\"omer, \emph{Simplifying Expressions Involving Radicals}, PhD thesis, Freie Universit\"at Berlin, 1993.
\access{https://cs.nyu.edu/exact/pap/rootBounds/sumOfSqrts/bloemerThesis.pdf}{Full dissertation PDF}. \status{T}
The most extensive single technical source in the consolidated core bibliography: independence, radical representations, denesting, and arithmetic/bit-size complexity. The thesis is not a second independent discovery of every theorem already published in the conference papers.

\bibitem{bloemer00} J. Bl\"omer, ``Denesting by bounded degree radicals,'' \emph{Algorithmica} \textbf{28} (2000), 2--15. \doi{10.1007/s004530010028}. Conference version: \emph{ESA 1997}, LNCS \textbf{1284}, 53--63, \doi{10.1007/3-540-63397-9_5}. \status{A}
Bounded allowed indices and optimal or near-optimal depth. The polynomial running-time parameter is the splitting-field description size, not automatically the original radical-expression length.

\bibitem{bloemer98} J. Bl\"omer, ``A probabilistic zero-test for expressions involving roots of rational numbers,'' \emph{ESA 1998}, LNCS \textbf{1461}, 151--162. \status{A}
Randomized exact-algebraic testing with a specified radical input representation. A separate contribution from the 1991 sums-of-radicals paper.

\bibitem{horng99} G. Horng and M.-D.~A. Huang, ``Solving polynomials by radicals with roots of unity in minimum depth,'' \emph{Mathematics of Computation} \textbf{68}(226) (1999), 881--885. \doi{10.1090/S0025-5718-99-01060-1}. \status{T}
A concise source for minimum-depth polynomial solving with roots of unity included explicitly in the model. Corrected terminal page: 885.

\bibitem{cotner} C.~F. Cotner, \emph{The Nesting Depths of Radical Expression}, PhD thesis, University of California, Berkeley, 1995.
\access{https://math.berkeley.edu/publications/nesting-depths-radical-expression}{Departmental dissertation record}. \status{A}
The singular ``Expression'' follows the institutional record. H.~W. Lenstra, Jr.\ is the advisor, not a second dissertation author. Included as a verified thesis lead; detailed results were not audited.

\bibitem{landau94} S. Landau, ``How to tangle with a nested radical,'' \emph{The Mathematical Intelligencer} \textbf{16}(2) (1994), 49--55. \doi{10.1007/BF03024284}. \status{T}
The most useful short orientation to the mathematical and algorithmic landscape before the general technical papers.

\bibitem{landau98} S. Landau, ``$\sqrt2+\sqrt3$: four different views,'' \emph{The Mathematical Intelligencer} \textbf{20} (1998), 55--60. \doi{10.1007/BF03025229}.
\access{https://privacyink.org/pdf/Sqrt2+Sqrt3.pdf}{Author-hosted PDF}. \status{T}
A focused exposition of one algebraic number from several perspectives; useful for separating representation, degree, and computation.

\bibitem{jeffrey} D.~J. Jeffrey and A.~D. Rich, ``Simplifying square roots of square roots by denesting,'' in M.~J. Wester (ed.), \emph{Computer Algebra Systems: A Practical Guide}, Wiley, 1999, 61--72. ISBN 0-471-98353-5.
\access{https://cybertester.com/data/denest.pdf}{Paper PDF}; \access{https://math.unm.edu/~wester/cas/book/contents.html}{editor's book contents}. \status{T}
Recursive practical denesting and complexity control. One of the two references explicitly cited by SymPy's square-root denester.

\bibitem{lenstra06} H.~W. Lenstra, Jr., \emph{Entangled Radicals}, AMS Colloquium Lectures, 2006.
\access{https://www.math.leidenuniv.nl/~hwl/papers/rad.pdf}{Lecture notes}. \status{B}
Background on interactions among radical extensions and cyclotomic structure. A structural companion, not a documented production denester.

\bibitem{perucca} A. Perucca, ``Linear relations among radicals,'' arXiv:2511.02498, 2025.
\access{https://arxiv.org/abs/2511.02498}{arXiv record}. \status{A}
A recent addition concerning the relation between radical-group indices and field-extension degrees, and the role of roots of unity in linear dependence. Relevant to basis construction.

\bibgroup{II. Ramanujan, special families, and higher indices}

\bibitem{ramanujan} S. Ramanujan, \emph{Collected Papers of Srinivasa Ramanujan}, G.~H. Hardy, P.~V. Seshu Aiyar, and B.~M. Wilson (eds.), Cambridge University Press, 1927. \status{B}
Primary historical collection. Individual finite identities and infinite-radical questions should be distinguished when following later references back to their origin.

\bibitem{berndtIV} B.~C. Berndt, \emph{Ramanujan's Notebooks, Part IV}, Springer, 1994. \status{B}
Principal book-length source for many algebraic and radical identities, with proofs and historical context supplied by the editor.

\bibitem{berndtV} B.~C. Berndt, \emph{Ramanujan's Notebooks, Part V}, Springer, 1998. \status{B}
Companion volume for the broader notebook context, units, and related identities. Not a general-purpose CAS algorithm reference.

\bibitem{berndt97} B.~C. Berndt, H.~H. Chan, and L.-C. Zhang, ``Ramanujan's association with radicals in India,'' \emph{American Mathematical Monthly} \textbf{104}(10) (1997), 905--911. \doi{10.1080/00029890.1997.11990738}.
\access{https://faculty.math.illinois.edu/~berndt/publications.html}{Author's publication directory}. \status{A}
Historical and mathematical context for Ramanujan's radical identities; a good entry point before the classification literature.

\bibitem{berndt98} B.~C. Berndt, H.~H. Chan, and L.-C. Zhang, ``Radicals and units in Ramanujan's work,'' \emph{Acta Arithmetica} \textbf{87}(2) (1998), 145--158.
\access{https://eudml.org/doc/207210}{EuDML record and access}. \status{A}
Explains the number-theoretic role of units in the identities. Distinct from the preceding historical article.

\bibitem{shanks} D. Shanks, ``Incredible identities,'' \emph{Fibonacci Quarterly} \textbf{12} (1974), 271, 280. \status{B}
A classical source of striking radical identities and a useful starting point for the subsequent corrections and revisitations.

\bibitem{spohn} W.~G. Spohn, Jr., ``Letter to the Editor,'' \emph{Fibonacci Quarterly} \textbf{14} (1976), 12. \status{B}
A corrective historical companion to the incredible-identities discussion; retained to avoid presenting that literature without its follow-up.

\bibitem{bumby} R.~T. Bumby, ``Incredible identities revisited,'' \emph{Fibonacci Quarterly} \textbf{25}(1) (1987), 62--64. \doi{10.1080/00150517.1987.12429728}.
\access{https://www.fq.math.ca/25-1.html}{Original volume index}. \status{A}
The issue year is 1987. A 2024 digitization/publication timestamp should not be mistaken for a recent theorem.

\bibitem{honsbeek99} M. Honsbeek, \emph{Denesting Certain Nested Radicals of Depth Two}, Report 9916, University of Nijmegen, 1999.
\access{https://repository.ubn.ru.nl/handle/2066/18722}{Institutional repository record}. \status{A}
The early report on the square root of a sum of cube roots. The repository identifier must not be relabelled as the later dissertation.

\bibitem{honsbeek05} M. Honsbeek, \emph{Radical Extensions and Galois Groups}, PhD thesis, Radboud University Nijmegen, 2005.
\access{https://www.math.ru.nl/~bosma/students/honsbeek/M_Honsbeek_thesis.pdf}{Full dissertation PDF}. \status{T}
Field-theoretic treatment of radical extensions, including the depth-two family and its denesting criterion. A substantial companion to the shorter 1999 report.

\bibitem{sury} B. Sury, \emph{Ramanujan's Route to Roots of Roots}, lecture notes, Indian Statistical Institute, Bangalore, for a talk at IIT Madras.
\access{https://www.isibang.ac.in/~sury/ramanujanday.pdf}{Institutional PDF}. \status{T}
Accessible treatment of the Galois-theoretic criterion, rational quartic test, and explicit formulas. The sign of $\beta$ in one printed example needs care; this report uses $(28,-27)$.

\bibitem{osler} T.~J. Osler, ``Cardan polynomials and the reduction of radicals,'' \emph{Mathematics Magazine} \textbf{74}(1) (2001), 26--32. \status{A}
Trace-type recurrence identities that explain some reductions of higher roots. Useful algebraic background for the Cardan route visible in practical denesters.

\bibitem{witula} R. Witu\l a and D. S\l ota, ``Cardano's formula, square roots, Chebyshev polynomials and radicals,'' \emph{Journal of Mathematical Analysis and Applications} \textbf{363} (2010), 639--647. \status{B}
Connects cubic formulas, Chebyshev-type identities, and radical expressions. Included as a related algebraic construction source, not a universal denesting theorem.

\bibitem{shevelev} V. Shevelev, ``On Ramanujan cubic polynomials,'' arXiv:0711.3420, 2007.
\access{https://arxiv.org/abs/0711.3420}{arXiv record}. \status{A}
Ramanujan cubic polynomials and associated radical identities. Useful before the cyclic-cubic and Shanks-polynomial connection in the next entry.

\bibitem{barbero} S. Barbero, U. Cerruti, N. Murru, and M. Abrate, ``Identities involving zeros of Ramanujan and Shanks cubic polynomials,'' \emph{Journal of Integer Sequences} \textbf{16} (2013), Article 13.8.1.
\access{https://cs.uwaterloo.ca/journals/JIS/VOL16/Barbero/barbero11.html}{Journal page and full text}; arXiv:1401.1474. \status{A}
Connects the two cubic families, cyclic cubic fields, Gaussian periods, trigonometric root formulas, and resulting identities. The 2014 arXiv posting is not the journal publication year.

\bibitem{krepkii} I.~A. Krepkii and K.~I. Pimenov, ``Cube root Ramanujan formulas and elementary Galois theory,'' \emph{Vestnik St. Petersburg University: Mathematics} \textbf{48} (2015), 214--223. \doi{10.3103/S106345411504007X}. \status{A}
Normal closures of pure cubic extensions over cyclic cubic fields, Hilbert's Theorem~90, and a prime-degree extension of the method. Online publication was in March 2016; the issue year is 2015.

\bibitem{antipov} M.~A. Antipov and K.~I. Pimenov, ``Ramanujan denesting formulas for cubic radicals,'' \emph{Vestnik St. Petersburg University: Mathematics} \textbf{53}(2) (2020), 115--121. \doi{10.1134/S1063454120020028}. \status{A}
A classification in a specified pure-cubic setting, including a three-summand bound and a norm condition. Its qualified uniqueness statement should not be generalized beyond that setting.

\bibitem{dokmeci} K. Dokmeci, \emph{Theorems on Field Extensions and Radical Denesting}, MIT PRIMES report, 2017.
\access{https://math.mit.edu/research/highschool/primes/materials/2017/Dokmeci.pdf}{Institutional PDF}. \status{B}
Student research on field extensions and real radical denesting, useful for alternative elementary routes. Author's given name: Kaan. Included with its actual publication status.

\bibitem{gkioulekas} E. Gkioulekas, ``On the denesting of nested square roots,'' \emph{International Journal of Mathematical Education in Science and Technology} \textbf{48}(6) (2017), 942--953. \doi{10.1080/0020739X.2017.1290831}.
\access{https://faculty.utrgv.edu/eleftherios.gkioulekas/papers/024-denesting-square-roots.pdf}{Author-hosted PDF}; \access{https://scholarworks.utrgv.edu/mss_fac/243/}{institutional record}. \status{T}
A clear elementary presentation of direct and indirect square-root denesting, useful alongside BFHT.

\bibitem{campbell} G.~B. Campbell and A. Zujev, ``Variations on Ramanujan's nested radicals,'' arXiv:1511.06865, 2015.
\access{https://arxiv.org/html/1511.06865v1}{Full HTML text}. \status{T}
Finite nested-radical identities, including cubic and quartic constructions. Not classified here as primarily an infinite-radical convergence paper.

\bibitem{girstmair} K. Girstmair, ``Reducing radicals in the spirit of Euclid,'' \emph{Journal of Symbolic Computation} \textbf{104} (2021), 356--365. \doi{10.1016/j.jsc.2020.07.019}; arXiv:2004.06039.
\access{https://arxiv.org/abs/2004.06039}{Preprint record}. \status{T}
Odd-index radical constructions viewed as algebraic-degree reduction. The 2020 preprint and 2021 journal paper are grouped; degree reduction is not automatically radical-depth reduction.

\bibitem{osipov21} N.~N. Osipov and A.~A. Kytmanov, ``Simplification of nested real radicals revisited,'' in \emph{Computer Algebra in Scientific Computing, CASC 2021}, LNCS \textbf{12865}, 293--313. \doi{10.1007/978-3-030-85165-1_17}. \status{A}
A detailed theory and algorithm for specified triply nested real radicals over $\Q$, with nonsimplifiable examples. Scope here is based on the publisher's abstract, not an invented unrestricted theorem.

\bibitem{cavallo} A. Cavallo, ``Denesting cubic radicals,'' arXiv:2403.04776, version 2, 28 September 2024.
\access{https://arxiv.org/html/2403.04776v2}{Full HTML text}; \access{https://sites.google.com/view/albertocavallomath/misc}{author's program page, as linked by the paper}. \status{T}
The quadratic-field cubic criterion is useful and is proved in rescaled form in Theorem~\ref{thm:cubic}. The preprint's integer-factorization inference is not accepted; the linked program's contents were not independently tested.

\bibitem{tonien} J. Tonien, ``Nested fifth root radical identities from elliptic curves,'' \emph{Journal of Computational Algebra} \textbf{13--14} (2025), Article 100032. \doi{10.1016/j.jaca.2025.100032}.
\access{https://www.sciencedirect.com/science/article/pii/S2772827725000038}{Publisher record}. \status{A}
Elliptic-curve constructions of fifth-root identities, with SageMath computational material described in the publication record. A higher-index family source, not a general denester.

\bibitem{osipov26} N.~N. Osipov, ``On Ramanujan's identities with nested cubic radicals,'' \emph{Programmirovanie} (Russian edition), 2026, no.~2, 71--77. \doi{10.7868/S3034584726020081}.
\access{https://journals.rcsi.science/0132-3474/article/view/446083}{Publisher page with English title and abstract}. \status{A}
A 2026 addition on specified cubic identity families and dependence on the base field. Full text required authentication; the damaged formulas in the accessible abstract are not reconstructed into unsupported theorem statements.

\bibgroup{III. Identity testing and exact-computation complexity}

\bibitem{rit} N. Balaji, K. Nosan, M. Shirmohammadi, and J. Worrell, ``Identity testing for radical expressions,'' \emph{LICS 2022}, Article 8, 1--11. \doi{10.1145/3531130.3533331}; arXiv:2202.07961, version 4, 16 October 2024.
\access{https://arxiv.org/abs/2202.07961}{Revised preprint record}. \status{A}
RIT for polynomial circuits evaluated at real radical constants; coNP under GRH, with stronger statements for the prime-square-root restriction. The circuit model and conditional hypotheses are essential.

\bibitem{gaillard} L. Gaillard and G. Jindal, ``On the order of power series and the sum of square roots problem,'' \emph{ISSAC 2023}. \doi{10.1145/3597066.3597079}; arXiv:2304.13605.
\access{https://arxiv.org/abs/2304.13605}{Preprint and abstract}. \status{A}
Wronskian/order methods and a succinct-input generalization of the equality problem. Not a general solution of square-root-sum sign comparison.

\bibitem{allender} E. Allender, P. B\"urgisser, J. Kjeldgaard-Pedersen, and P.~B. Miltersen, ``On the complexity of numerical analysis,'' \emph{SIAM Journal on Computing} \textbf{38}(5) (2009), 1987--2006. \status{A}
Foundational complexity background for exact numerical decisions, including square-root-sum problems. Included without an unsupported assertion that every bound remains the best known as of the source-check date.

\bibgroup{IV. Books and broader mathematical background}

\bibitem{zippelbook} R. Zippel, \emph{Effective Polynomial Computation}, Kluwer Academic Publishers, 1993. \status{B}
Polynomial and symbolic-computation background from a principal contributor to radical simplification. Useful for the structural and representation issues surrounding the papers.

\bibitem{geddes} K.~O. Geddes, S.~R. Czapor, and G. Labahn, \emph{Algorithms for Computer Algebra}, Kluwer Academic Publishers, 1992. \status{B}
Core algorithms for exact symbolic computation, including algebraic extensions and polynomial arithmetic needed by denesters.

\bibitem{gathen} J. von zur Gathen and J. Gerhard, \emph{Modern Computer Algebra}, 3rd ed., Cambridge University Press, 2013. \status{B}
Algorithmic infrastructure and complexity analysis. Not cited as containing a complete modern survey of minimum-depth denesting.

\bibitem{davenport} J.~H. Davenport, Y. Siret, and \'E. Tournier, \emph{Computer Algebra}, Academic Press, 1988. \status{B}
General CAS and algebraic-simplification background; English edition of the work originally published in French.

\bibitem{cohen} H. Cohen, \emph{A Course in Computational Algebraic Number Theory}, Graduate Texts in Mathematics \textbf{138}, Springer, 1993. \status{B}
Number-field arithmetic, factorization, and representations supporting rigorous algebraic implementations.

\bibitem{cox} D.~A. Cox, \emph{Galois Theory}, 2nd ed., Wiley, 2012. \status{B}
Background on radical extensions, solvable groups, roots of unity, and the distinction between a polynomial and its splitting field.

\bibitem{newman} S.~C. Newman, \emph{A Classical Introduction to Galois Theory}, Wiley, 2012. ISBN 978-1-118-33681-6. \status{B}
An additional route into the field-theoretic language used throughout the denesting literature.

\bibitem{euclid} Euclid, \emph{Elements}, Book X. \status{B}
Historical background on irrational magnitudes and quadratic constructions. Modern generalizations and degree-reduction results should be attributed separately to their modern authors.

\bibitem{piagm} J.~M. Borwein and P.~B. Borwein, \emph{Pi and the AGM: A Study in Analytic Number Theory and Computational Complexity}, Wiley, 1987. \status{B}
Adjacent background on algebraic transformations, infinite products, and analytic computation. Not a finite-denesting algorithm source.

\bibgroup{V. Infinite radicals and related implicit equations}

\bibitem{herschfeld} A. Herschfeld, ``On infinite radicals,'' \emph{American Mathematical Monthly} \textbf{42} (1935), 419--429. \status{A}
Classical convergence criterion for nonnegative nested square roots. The bounded sequence is $a_n^{2^{-n}}$, not necessarily $a_n$ itself.

\bibitem{borwein} J.~M. Borwein and G. de Barra, ``Nested radicals,'' \emph{American Mathematical Monthly} \textbf{98} (1991), 735--739. \status{A}
A later analytic discussion of nested radicals and convergence, retained separately from finite algebraic denesting.

\bibitem{elia} M. Elia, ``Observations on some algebraic equations associated with Ramanujan's work,'' in \emph{Number Theory and Discrete Mathematics}, Birkh\"auser, 2002, 101--111. \doi{10.1007/978-3-0348-8223-1_11}. \status{A}
Implicit algebraic equations produced by iterated radical patterns. The correct author is Michele Elia; the attribution to Berndt and Andrews in a supplied report is incorrect.

\bibitem{mitra} N. Mitra, ``Generalization of Ramanujan's famous nested radicals to the $n$th root and their evaluation,'' arXiv:2404.04051, 2024.
\access{https://arxiv.org/abs/2404.04051}{Preprint record}. \status{A}
Generalizations of the famous infinite-radical evaluations. Included as a related analytic subject, not as a decision algorithm for finite denesting.

\bibitem{rahal} S. Rahal, \emph{Introduction to Nested Radicals}, bachelor's thesis, Stockholm University, 2017, report 2017:17.
\access{https://kurser.math.su.se/pluginfile.php/16103/mod_folder/content/0/2017/2017_examensarbeten_matematik.pdf?forcedownload=1}{Institutional 2017 thesis register}. \status{A}
Introductory thesis lead on nested radicals. Institutional metadata checked; detailed results are not used as the basis of a new theorem claim.

\bibitem{infiniteblogs} Ancillary infinite-radical exposition: ``Ramanujan's nested radical problem,'' \emph{Cantor's Paradise},
\access{https://www.cantorsparadise.com/ramanujans-nested-radical-problem-f14c3060e495}{article}; A.~G. Smith, ``Calculus of infinite nested radicals,''
\access{https://medium.com/@ArchieGSmith/calculus-of-infinite-nested-radicals-1ab3e1d26774}{Medium article}; and
\access{https://math.stackexchange.com/questions/4049946/the-infinitely-nested-radicals-problem-and-ramanujans-wondrous-formula}{Math StackExchange discussion 4049946}. \status{B}
Informal leads retained from the supplied material. They are not used to support the finite-denesting criteria or complexity claims.

\bibgroup{VI. Official software documentation, source, and development discussions}

\bibitem{sympydocs} SymPy developers, \emph{Simplify documentation}, entries for \texttt{sqrtdenest}, \texttt{nthroot}, \texttt{radsimp}, and \texttt{powdenest}.
\access{https://docs.sympy.org/latest/modules/simplify/simplify.html}{Official documentation}. \status{D}
Public contracts and examples. The real $n$th-root routine is a significant addition to the coverage in the supplied reports. The local verification run used version 1.14.0.

\bibitem{sympysrc} SymPy developers, \texttt{sympy/simplify/sqrtdenest.py}, version 1.14.0, and the associated test suite.
\access{https://github.com/sympy/sympy/blob/1.14.0/sympy/simplify/sqrtdenest.py}{Versioned source};
\access{https://github.com/sympy/sympy/blob/1.14.0/sympy/simplify/tests/test_sqrtdenest.py}{versioned tests}. \status{D}
Implementation-level reading beside BFHT and Jeffrey--Rich. Private helper names are navigation aids and may change between releases.

\bibitem{sympyissue} SymPy development discussions:
\access{https://github.com/sympy/sympy/issues/2415}{issue 2415} and
\access{https://github.com/sympy/sympy/issues/27082}{issue 27082, ``simplify roots with base having surds''}. \status{D}
Regression and feature-discussion leads. They are not evidence that the reported behavior persists in every current API; \texttt{nthroot} handles examples that a square-root-only route does not.

\bibitem{raddenestsrc} G. Schintgen, \emph{raddenest}, Maxima package.
\access{https://github.com/gschintgen/raddenest}{Repository};
\access{https://raw.githubusercontent.com/gschintgen/raddenest/master/raddenest.mac}{implementation source}. \status{D}
A SymPy-derived square-root denester with original higher-root extensions, including Cardan and Honsbeek-style routines. Source comments and branch settings should be read before translating examples involving negative radicands.

\bibitem{maximamanual} Maxima developers, \emph{Maxima Manual}.
\access{https://maxima.sourceforge.io/docs/manual/maxima_singlepage.html}{Official manual}. \status{B}
Reference directory for \texttt{radcan}, radical expansion, and installed package interfaces. Legacy \texttt{sqdnst}/\texttt{sqrtdenest} release-history claims were not independently settled here.

\bibitem{maximathread} G. Schintgen and participants, denesting and \texttt{raddenest} discussions, \emph{maxima-discuss}, 2017.
\access{https://sourceforge.net/p/maxima/mailman/maxima-discuss/thread/o6lan6\%24h4r\%241\%40blaine.gmane.org/}{SourceForge thread};
\access{https://maxima-discuss.narkive.com/rz7bUs4t/denesting}{discussion mirror}. \status{B}
Historical development context. The package source, not a historical message alone, is the main evidence for capabilities described in this report.

\bibitem{wolframroot} Wolfram Research, \emph{RootReduce}, Wolfram Language documentation.
\access{https://reference.wolfram.com/language/ref/RootReduce.html}{Official reference}. \status{D}
Canonical exact algebraic-number reduction. Its output language may be \texttt{Root}, so it is not in itself a minimum-depth radical-expression interface.

\bibitem{wolframradical} Wolfram Research, \emph{ToRadicals}, Wolfram Language documentation.
\access{https://reference.wolfram.com/language/ref/ToRadicals.html}{Official reference}. \status{D}
Low-degree algebraic-root conversion, higher-degree limitations, and parameter qualifications. Conversion from an algebraic number to radicals is distinct from optimizing a pre-existing radical expression.

\bibitem{wolframrd} S. Banerjee, \emph{RadicalDenest}, Wolfram Function Repository, version 2.1.0, 29 April 2024.
\access{https://resources.wolframcloud.com/FunctionRepository/resources/RadicalDenest}{Function page}. \status{D}
A heuristic, time-constrained denesting resource function; acknowledges Corey Ziegler, Bill Gosper, and Daniel Lichtblau. Not a built-in kernel guarantee of complete minimum-depth denesting.

\bibitem{wolframcommunity} Wolfram Community, ``Heuristic package to denest radicals'' and ``Radical Denest: an ancient difficult task in symbolic computation.''
\access{https://community.wolfram.com/groups/-/m/t/980264}{Discussion 980264};
\access{https://community.wolfram.com/groups/-/m/t/2120629}{discussion 2120629}. \status{D}
Development discussion and examples underlying the contributed heuristic tools. Useful for provenance and design trade-offs, not as a replacement for formal completeness proofs.

\bibitem{maplerad} Maplesoft, \emph{radnormal}, Maple help.
\access{https://www.maplesoft.com/support/help/Maple/view.aspx?path=radnormal}{Official documentation}. \status{D}
Describes exact radical normalization, zero detection, principal roots, and the indexed-\texttt{RootOf}/field-basis pipeline. Denesting is explicitly only sometimes achieved.

\bibitem{maplesimp} Maplesoft, \emph{simplify/radical}, Maple help.
\access{https://www.maplesoft.com/support/help/Maple/view.aspx?path=simplify\%2Fradical}{Official documentation}. \status{D}
Associated radical simplification and assumptions. Read together with \texttt{radnormal} instead of attributing an unspecified universal denester to \texttt{simplify}.

\bibitem{sageaa} SageMath developers, \emph{Field of algebraic numbers}, \texttt{sage.rings.qqbar} reference, especially \texttt{radical\_expression()}.
\access{https://doc.sagemath.org/html/en/reference/number_fields/sage/rings/qqbar.html}{Official reference}. \status{D}
Exact real/complex algebraic-number representations and attempted radical conversion. An unchanged exact algebraic object is not a numerical approximation or impossibility certificate.

\bibitem{sagesymbolic} SageMath developers, \emph{Symbolic expressions}, especially \texttt{canonicalize\_radical()}.
\access{https://doc.sagemath.org/html/en/reference/calculus/sage/symbolic/expression.html}{Official reference}. \status{D}
Symbolic radical normalization with explicit branch/assumption cautions. Distinct from the algebraic-field conversion method in the preceding entry.

\bibitem{pari} PARI/GP developers, \emph{General number fields}, official reference.
\access{https://pari.math.u-bordeaux.fr/dochtml/html-stable/General_number_fields.html}{Number-field documentation}. \status{D}
Subfields, embeddings, factorization, and exact field arithmetic are building blocks for a denester. They do not automatically reconstruct a minimum-depth expression for a selected element.

\bibitem{magma} Magma Computational Algebra System, \emph{Handbook} and number-field release notes.
\access{https://magma.maths.usyd.edu.au/magma/handbook/}{Handbook};
\access{https://magma.maths.usyd.edu.au/magma/releasenotes/pdf/relv210.pdf}{Version 2.10 release notes}. \status{D}
Optimized field representations and related operations concern defining data. The record is not used to assert the absence of every radical simplification facility.

\bibitem{fricas} FriCAS developers, \emph{RadicalSolvePackage} API; related Axiom documentation.
\access{https://fricas.github.io/api/RadicalSolvePackage.html}{Official FriCAS API}. \status{D}
Radical polynomial-solving interfaces. Newer \texttt{rsimp}/\texttt{recursiveRootSimplification} claims from the supplied material remain unverified at a precise release/API level.

\bibitem{honda} Y. Honda, \emph{GaloisGroupSolver}, Maxima package.
\access{https://github.com/YasuakiHonda/GaloisGroupSolver}{Author's repository}. \status{B}
Adjacent solvable-polynomial/Galois-group implementation lead. Not tested here and not asserted to be a minimum-depth denester.

\bibgroup{VII. Encyclopedias, elementary pages, and question-and-answer resources}

\bibitem{wiki} Wikipedia contributors, \emph{Nested radical}, especially ``Denesting.''
\access{https://en.wikipedia.org/wiki/Nested_radical\#Denesting}{Article section}. \status{B}
The original prompt's starting point; useful for orientation and bibliography. Core results in this report are supported by primary papers or direct proofs rather than this page alone.

\bibitem{mathworld} E.~W. Weisstein, \emph{Nested Radical}, MathWorld.
\access{https://mathworld.wolfram.com/NestedRadical.html}{Encyclopedia entry}. \status{B}
A compact source of identities and literature leads. Check domain and denesting-model restrictions in the original references.

\bibitem{brown} S. Brown, \emph{Denesting Radicals (or Unnesting Radicals)}.
\access{https://brownmath.com/alge/nestrad.htm}{Author's tutorial}. \status{T}
Worked elementary square-root examples. A convenient first algebraic introduction, not a general field-theoretic classification.

\bibitem{mo} MathOverflow, \emph{Relevance of Landau's algorithm for denesting radicals}, question 319874; related discussions 410388 and 344478.
\access{https://mathoverflow.net/questions/319874/relevance-of-landaus-algorithm-for-denesting-radicals}{Question 319874};
\access{https://mathoverflow.net/questions/410388}{question 410388};
\access{https://mathoverflow.net/questions/344478}{question 344478}. \status{B}
Research-level discussion leads on general algorithms, multiple square roots, and rationality/representation. Not treated as independent proof of a theorem already cited from the literature.

\bibitem{mse} Math StackExchange, discussions of elementary radical denesting and Ramanujan cube-root identities:
\access{https://math.stackexchange.com/questions/4680}{question 4680},
\access{https://math.stackexchange.com/questions/16331}{question 16331}, and
\access{https://math.stackexchange.com/questions/871639}{question 871639}. \status{B}
Concrete question identifiers replace vague unverified titles from the supplied reports. Useful coefficient-matching, field-theoretic, and exact-identity explanations; informal rather than a substitute for primary-source verification.

\end{thebibliography}
\end{document}
